\section{Model Design}
\label{sec:method}

\begin{figure*}[h]
    \centering
    \includegraphics[width=\linewidth, trim=10 140 40 120, clip]{figures/pipeline_overview.pdf}
    \vspace{-10pt}
    \caption{
        Overview of the \modelname{} framework.
    }
    \label{fig:pipeline}
\end{figure*}






The overview framework of \modelname{} is shown in Fig. \ref{fig:pipeline}.
The core is our HY-Motion DiT model, which accepts a text prompt and an expected duration of the motion as inputs and generates a clip of 3D human motion as outputs.
We first detail our motion representation in \cref{subsec:motion_rep}.
Then the HY-Motion DiT model is described in \cref{subsec:backbone}.
Besides the core model, we train an independent LLM (\cref{subsec:pipeline_enhancement}) to analyze user prompts and predict the expected duration of the motion clips, and additionally, convert user prompts into better-structured prompts.






\subsection{Motion Representation}
\label{subsec:motion_rep}



We employ the skeleton definition of SMPL-H~\citep{Romero2017} (22 joints without hands).
Formally, we denote a motion sequence as $\bm{x}=\{\bm{f}_1,\bm{f}_2,\dots,\bm{f}_N\}$, where each frame is represented as a vector $\bm{f}\in\mathbb{R}^{201}$, comprising the global root translation $\bm{t}\in\mathbb{R}^3$, the global body orientation $\bm{r}\in\mathbb{R}^6$, the local joint rotations $\bm{j}^r\in\mathbb{R}^{21\times6}$, and the local joint positions $\bm{j}^p\in\mathbb{R}^{22\times3}$.
All rotational parameters adhere to the continuous 6D representation~\citep{Zhou2019}.
Noticeably, our representation is similar to DART~\citep{DartControl}, where the rotation representation is compatible with typical animation workflows, which differs from the commonly used HumanML3D~\citep{Guo2022} representation.
We further remove the explicit temporal derivatives (velocities) and foot contact labels, as we observed faster training convergence.






\subsection{HY-Motion DiT}
\label{subsec:backbone}

\begin{figure*}[h]
    \centering
    \includegraphics[width=\linewidth]{figures/backbone_mmdit.pdf}
    \vspace{-10pt}
    \caption{
        Model architecture of our HY-Motion DiT.
    }
    \label{fig:backbone_mmdit}
\end{figure*}

\paragraph{Model Architecture.}
Similar to HunyuanVideo~\cite{Kong2024}, we employ a hybrid Transformer that combines dual-stream and single-stream processing to model the joint distribution of motion and text.
The network begins with dual-stream blocks, where motion latents and text tokens are processed via independent QKV projections and MLPs.
Crucially, they interact through a joint attention mechanism, allowing motion features to query semantic cues from the text while preserving their distinct, modality-specific representations.
These streams subsequently merge in the single-stream blocks, where motion and text tokens are concatenated into a unified sequence.
Here, we employ parallel spatial and channel attention modules to facilitate deep multimodal fusion and information exchange.


\paragraph{Text Encoders.}
We leverage a hierarchical dual-conditioning strategy to integrate both fine-grained and global text guidance.
Firstly, we employ Qwen3-8B~\citep{Yang2025} to extract rich, token-wise semantic embeddings.
Addressing the limitation that LLMs typically utilize causal attention—which restricts the context for non-autoregressive generation, we follow HunyuanVideo~\citep{Kong2024} by adopting a Bidirectional Token Refiner.
This module transforms the causal LLM features into bidirectional representations before injecting them into the dual-stream blocks, thereby providing holistically contextualized guidance for the diffusion model.
Complementarily, we utilize CLIP-L~\citep{Radford2021} to extract a global text embedding.
This embedding is concatenated with the timestep embedding and injected via a separated AdaLN~\citep{Peebles2023} mechanism, where layer-specific modulation parameters are learned to adaptively regulate feature statistics throughout the network.


\paragraph{Attention Mechanism and Positional Encoding.}
We devise a composite Attention Mechanism that incorporates specific masking strategies for both cross-modal interaction and temporal modeling.
First, to regulate multimodal information flow, we enforce an asymmetric attention mask.
While motion tokens attend globally to the text sequence to extract semantic cues, text tokens are explicitly masked from the motion latents.
This unidirectional constraint prevents the diffusion noise inherent in the motion states from propagating back to the text embeddings, thereby preserving the integrity of the semantic conditioning.
Second, for temporal modeling within the motion branch, we implement a narrow band mask strategy.
Predicated on the hypothesis that kinematic dynamics are governed primarily by local continuity, we restrict attention to a sliding window of $121$ frames under $30$ fps.
This imposes a locality inductive bias that decomposes long sequences into coherent substructures, effectively focusing the model on complex local dynamics while ensuring linear computational complexity.
To spatially ground these attention interactions, we optimize the Positional Encoding by adopting full Rotary Positional Embeddings (RoPE) \cite{su2024roformer}.
Diverging from disjoint encoding schemes, we concatenate text and motion embeddings into a single sequence prior to encoding.
By applying RoPE across this unified sequence, we establish a continuous relative coordinate system, enabling the model to inherently resolve the positional correspondence between specific textual tokens and temporal motion frames.

\paragraph{Flow Matching Objective.}
We employ Flow Matching~\citep{Lipman2022} to construct a continuous probability path that bridges the standard Gaussian noise distribution and the complex motion data distribution.
We adopt the optimal transport path, defined as a linear interpolation $\bm{x}_t=(1-t)\bm{x}_0 + t\bm{x}_1$, which implies a constant target velocity~\citep{Liu2022}.
The training objective is to minimize the Mean Squared Error (MSE) between the predicted and ground-truth velocity:
\begin{equation}
    \mathcal{L}_{\text{FM}} = \mathbb{E}_{t, \bm{x}_0, \bm{x}_1} [ || \bm{v}_{\bm{\theta}}(\bm{x}_t,\bm{c},t) - \bm{v}_t ||^2_2 ],
    \label{eq:flow_matching}
\end{equation}
where $\bm{x}_1$ represents the clean motion data, $\bm{x}_0\sim\mathcal{N}(0, \bm{I})$ is the initial noise, and the target velocity is given by $\bm{v}_t=\bm{x}_1-\bm{x}_0$.
During inference, the generation process is formulated as an Ordinary Differential Equation (ODE): $d\bm{x}/d\bm{t}=\bm{v}_{\bm{\theta}}(\bm{x}_t,\bm{c},t)$.
Starting from random noise $\bm{x}_0$, we recover the clean motion $\bm{x}_1$ by numerically integrating this ODE along the predicted velocity field using an ODE solver (\eg, Euler).


\subsection{Duration Prediction \& Prompt Rewrite}
\label{subsec:pipeline_enhancement}

To accommodate diverse user inputs, we employ a dedicated LLM for both duration prediction and prompt rewrite. This module first predicts the target motion's duration from the user prompt. It then transforms the prompt into a structured format optimized for our DiT model. This design is motivated by the observation that LLMs possess inherent common-sense knowledge about the typical duration of human activities, enabling them to infer temporal length from textual descriptions. To further enhance its accuracy, we fine-tune the LLM on a dataset of ground-truth motion durations, which aligns its predictions with our training data distribution. Crucially, this fine-tuning preserves the model's ability to generalize to unseen motion descriptions.


\paragraph{Data Synthesis.}
Our dataset for fine-tuning the LLM consists of data triplets \{\textit{user prompt}, \textit{optimized prompt}, \textit{duration}\}. The ``optimized prompt'' and ``duration'' are the ground-truth text caption and duration of each motion clip in our motion dataset, respectively.
To simulate realistic ``user prompt'', we employed a powerful LLM (\eg, Gemini-2.5-Pro) to synthesize a dataset of user queries. Using the ground-truth text captions as a basis, we prompted the LLM to generate inputs that mimic the nature of casual human commands. These synthetic queries are intentionally diverse, encompassing informal language, a mix of English and Chinese, and varying levels of specificity, including queries that are deliberately brief and ambiguous.



\paragraph{Two-Stage Fine-Tuning.}
The duration prediction \& prompt rewrite module is fine-tuned from a Qwen3-30B-A3B~\citep{Yang2025} model with a two-stage training procedure.
\begin{itemize}
    \item \textit{Supervised Fine-Tuning (SFT):}
    The model is first fine-tuned on the constructed dataset with ground-truth supervision.
    It learns to sequentially generate ``optimized prompt'' and ``duration'' from the ``user prompt'', effectively grounding the temporal prediction in the semantic context of the rewritten text.

    \item \textit{Reinforcement Learning (RL):}
    To further mitigate hallucinations and enhance fidelity, we reinforce the model using Group Relative Policy Optimization (GRPO)~\citep{Shao2024}.
    We utilize a more powerful model (Qwen3-235B-A22B-Instruct-2507~\citep{Yang2025}) as a reward judge to evaluate multiple candidate outputs for each prompt.
    The reward function assesses two critical dimensions: semantic consistency, ensuring the rewrite remains faithful to the user's intent, and temporal plausibility, verifying that the predicted duration aligns with the physical complexity of the described action.
    By optimizing against the relative advantages of these candidates, GRPO steers the policy toward generating instructions that are both semantically precise and temporally coherent.
\end{itemize}





